% Section 08 - Dépannage et résolution de problèmes
\section{Dépannage et résolution de problèmes}

\subsection{Problèmes courants et solutions}

\subsubsection{L'application ne se lance pas}

\textbf{Symptômes :}
\begin{itemize}[leftmargin=*]
    \item Double-clic sur l'icône sans effet
    \item Fenêtre s'affiche brièvement puis disparaît
    \item Message d'erreur au démarrage
\end{itemize}

\textbf{Solutions :}
\begin{enumerate}[leftmargin=*]
    \item \textbf{Vérifier la configuration système}
    \begin{itemize}
        \item RAM disponible : minimum 4 Go
        \item Espace disque : minimum 500 Mo
        \item Système d'exploitation à jour
    \end{itemize}

    \item \textbf{Désinstaller et réinstaller}
    \begin{itemize}
        \item Désinstallez complètement l'application
        \item Téléchargez la dernière version
        \item Réinstallez
    \end{itemize}

    \item \textbf{Vérifier les permissions}
    \begin{itemize}
        \item Windows : Exécuter en tant qu'administrateur
        \item macOS : Autoriser dans Sécurité et confidentialité
        \item Linux : Vérifier les permissions d'exécution
    \end{itemize}
\end{enumerate}

\subsubsection{Erreur lors du chargement du fichier Excel}

\textbf{Message d'erreur :} "Impossible de lire le fichier Excel" ou "Colonnes manquantes"

\textbf{Solutions :}
\begin{enumerate}[leftmargin=*]
    \item \textbf{Vérifier le format du fichier}
    \begin{itemize}
        \item Format accepté : .xlsx ou .xls
        \item Pas de macros actives
        \item Fichier non protégé par mot de passe
    \end{itemize}

    \item \textbf{Vérifier les colonnes obligatoires}
    \begin{itemize}
        \item Pour relevés : MATRICULE, NOM, PRENOM, SEMESTRE, NIVEAU, FILIERE
        \item Pour attestations : Ajouter DATE\_NAISSANCE, LIEU\_NAISSANCE, DOMAINE, PARCOURS, MOYENNE\_GENERALE
        \item Respect de la casse (majuscules/minuscules)
    \end{itemize}

    \item \textbf{Fermer Excel}
    \begin{itemize}
        \item Le fichier ne doit pas être ouvert dans Excel
        \item Fermez Excel complètement
        \item Réessayez le chargement
    \end{itemize}

    \item \textbf{Vérifier l'encodage}
    \begin{itemize}
        \item Enregistrez le fichier avec encodage UTF-8
        \item Évitez les caractères spéciaux dans les noms de colonnes
    \end{itemize}
\end{enumerate}

\subsubsection{Erreur "Out of Memory" (OOM)}

\textbf{Message :} "JavaScript heap out of memory" ou "Mémoire insuffisante"

\textbf{Causes :}
\begin{itemize}[leftmargin=*]
    \item Génération de trop de documents simultanément
    \item Fichier Excel très volumineux
    \item Autres applications consommant beaucoup de RAM
\end{itemize}

\textbf{Solutions :}
\begin{enumerate}[leftmargin=*]
    \item \textbf{Réduire le nombre de documents}
    \begin{itemize}
        \item Limitez à 200 relevés par génération
        \item Pour plus, divisez en plusieurs lots
        \item Utilisez les filtres pour sélectionner un sous-ensemble
    \end{itemize}

    \item \textbf{Fermer les autres applications}
    \begin{itemize}
        \item Libérez de la RAM
        \item Fermez les navigateurs avec beaucoup d'onglets
        \item Redémarrez l'ordinateur si nécessaire
    \end{itemize}

    \item \textbf{Simplifier le fichier Excel}
    \begin{itemize}
        \item Supprimez les colonnes inutiles
        \item Retirez les formules complexes
        \item Conservez uniquement les données nécessaires
    \end{itemize}
\end{enumerate}

\subsubsection{Mapping des colonnes non sauvegardé}

\textbf{Symptôme :} À chaque chargement d'Excel, le mapping doit être refait

\textbf{Cause :} Changement des noms de colonnes entre les chargements

\textbf{Solution :}
\begin{itemize}[leftmargin=*]
    \item Gardez les mêmes noms de colonnes dans vos fichiers Excel
    \item Le mapping est basé sur le NOM de la colonne, pas sa position
    \item Si vous devez renommer, utilisez \bouton{Réinitialiser le mapping}
\end{itemize}

\subsubsection{QR codes non lisibles}

\textbf{Symptômes :}
\begin{itemize}[leftmargin=*]
    \item Le scanner ne détecte pas le QR code
    \item Message d'erreur lors du scan
    \item QR code trop petit ou trop flou
\end{itemize}

\textbf{Solutions :}
\begin{enumerate}[leftmargin=*]
    \item \textbf{Augmenter la taille}
    \begin{itemize}
        \item Passez de "Petit" à "Moyen" ou "Grand"
        \item Minimum recommandé : 80x80 pixels
    \end{itemize}

    \item \textbf{Augmenter la correction d'erreur}
    \begin{itemize}
        \item Passez de "L" ou "M" à "Q" ou "H"
        \item Plus de redondance = meilleure lisibilité
    \end{itemize}

    \item \textbf{Réduire le contenu}
    \begin{itemize}
        \item Limitez le nombre de colonnes dans le QR
        \item Trop de données = QR trop dense
        \item Maximum recommandé : 5-6 champs
    \end{itemize}

    \item \textbf{Vérifier la position}
    \begin{itemize}
        \item Le QR ne doit pas chevaucher du texte
        \item Fond blanc recommandé autour du QR
        \item Évitez les zones avec images de fond
    \end{itemize}
\end{enumerate}

\subsubsection{Moyennes calculées incorrectement}

\textbf{Symptôme :} Les moyennes affichées ne correspondent pas aux calculs manuels

\textbf{Vérifications :}
\begin{enumerate}[leftmargin=*]
    \item \textbf{Poids des ECs}
    \begin{itemize}
        \item Vérifiez les coefficients dans la configuration
        \item Formule : $\text{Moy UE} = \frac{\sum(\text{Note} \times \text{Poids})}{\sum \text{Poids}}$
    \end{itemize}

    \item \textbf{Crédits des UEs}
    \begin{itemize}
        \item Les crédits ECTS sont-ils corrects ?
        \item Formule : $\text{Moy Gén} = \frac{\sum(\text{Moy UE} \times \text{Crédits})}{\sum \text{Crédits}}$
    \end{itemize}

    \item \textbf{Notes manquantes}
    \begin{itemize}
        \item Les cellules vides sont ignorées dans le calcul
        \item Assurez-vous que toutes les notes sont présentes
    \end{itemize}

    \item \textbf{Base de notation}
    \begin{itemize}
        \item Vérifiez que la base définie (10, 20, 100) correspond aux notes
        \item L'application normalise automatiquement sur 20
    \end{itemize}
\end{enumerate}

\subsubsection{Attestations générées pour des étudiants ajournés}

\textbf{Symptôme :} Des étudiants avec moyenne < 10 reçoivent des attestations

\attention{Ceci ne devrait PAS se produire. L'application filtre automatiquement.}

\textbf{Vérifications :}
\begin{enumerate}[leftmargin=*]
    \item Vérifiez la colonne \texttt{MOYENNE\_GENERALE} dans l'Excel
    \item Les valeurs doivent être numériques (pas de texte)
    \item Format : 10.5, 12.75, etc. (point décimal, pas virgule)
    \item Si le problème persiste, contactez le support
\end{enumerate}

\subsubsection{Thème personnalisé non appliqué}

\textbf{Symptôme :} Le thème configuré n'apparaît pas sur les documents générés

\textbf{Solutions :}
\begin{enumerate}[leftmargin=*]
    \item Vérifiez que le thème est bien enregistré (bouton \bouton{Enregistrer le thème})
    \item Pour les relevés : le thème est lié à la configuration de classe
    \item Pour les attestations : le thème est indépendant
    \item Essayez de prévisualiser pour voir si le thème s'applique
    \item Exportez et réimportez le thème pour tester
\end{enumerate}

\subsection{Optimisation des performances}

\subsubsection{Accélérer la génération}

\begin{enumerate}[leftmargin=*]
    \item \textbf{Désactiver la compression ZIP} (si non nécessaire)
    \begin{itemize}
        \item Gain de temps : 20-30\%
        \item Fichiers plus volumineux
    \end{itemize}

    \item \textbf{Générer par lots de 100-200}
    \begin{itemize}
        \item Au lieu de 500 en une fois
        \item Évite les erreurs mémoire
        \item Permet de vérifier progressivement
    \end{itemize}

    \item \textbf{Simplifier les thèmes}
    \begin{itemize}
        \item Évitez trop d'images ou de polices complexes
        \item Utilisez des préréglages plutôt que des thèmes très personnalisés
    \end{itemize}

    \item \textbf{Utiliser un SSD}
    \begin{itemize}
        \item Plus rapide qu'un disque dur mécanique
        \item Améliore l'écriture des fichiers générés
    \end{itemize}
\end{enumerate}

\subsubsection{Réduire l'utilisation mémoire}

\begin{enumerate}[leftmargin=*]
    \item Fermez toutes les applications inutiles
    \item Utilisez le format "Fichiers individuels" au lieu de "ZIP" si possible
    \item Ne générez pas de QR codes si non nécessaire
    \item Redémarrez l'application entre deux générations massives
\end{enumerate}

\subsection{Journalisation et logs}

\subsubsection{Consulter les logs de l'application}

Les logs peuvent aider à identifier des problèmes :

\textbf{Windows :}
\begin{verbatim}
%APPDATA%\FMSP-Releves\logs\
\end{verbatim}

\textbf{macOS :}
\begin{verbatim}
~/Library/Logs/FMSP-Releves/
\end{verbatim}

\textbf{Linux :}
\begin{verbatim}
~/.config/FMSP-Releves/logs/
\end{verbatim}

\subsubsection{Console de développement}

Pour les utilisateurs avancés :

\begin{enumerate}[leftmargin=*]
    \item Appuyez sur \texttt{F12} ou \texttt{Ctrl+Shift+I} (Windows/Linux)
    \item Ou \texttt{Cmd+Option+I} (macOS)
    \item La console de développement s'ouvre
    \item Consultez l'onglet \textbf{Console} pour les messages
    \item Recherchez les messages en rouge (erreurs)
\end{enumerate}

\subsection{Réinitialisation de l'application}

Si les problèmes persistent, réinitialisez l'application :

\begin{tcolorbox}[colback=red!5!white,colframe=red!75!black,title=ATTENTION : Sauvegardez vos données]
La réinitialisation supprime :
\begin{itemize}[leftmargin=*]
    \item Toutes les configurations de classes
    \item Tous les thèmes personnalisés
    \item L'historique des générations
    \item Les paramètres de l'établissement
\end{itemize}

Exportez vos configurations et thèmes avant de continuer !
\end{tcolorbox}

\textbf{Procédure :}
\begin{enumerate}[leftmargin=*]
    \item Menu \menu{Paramètres} > \textbf{Avancé}
    \item Cliquez sur \bouton{Réinitialiser l'application}
    \item Confirmez l'action
    \item L'application redémarre avec les paramètres par défaut
\end{enumerate}

\subsection{Contacter le support}

Si aucune solution ne fonctionne :

\begin{enumerate}[leftmargin=*]
    \item Menu \menu{Aide} > \bouton{Rapporter un problème}
    \item Remplissez le formulaire :
    \begin{itemize}
        \item Description du problème
        \item Étapes pour reproduire
        \item Captures d'écran si possible
    \end{itemize}
    \item Joignez les logs récents
    \item Cliquez sur \bouton{Envoyer}
\end{enumerate}

\textbf{Informations utiles à fournir :}
\begin{itemize}[leftmargin=*]
    \item Version de l'application (Menu \menu{Aide} > \menu{À propos})
    \item Système d'exploitation et version
    \item Nombre d'étudiants traités
    \item Type de document généré
    \item Message d'erreur exact (copier-coller)
\end{itemize}

\newpage
